%% ======================================================================
\section{Self-bounding algorithms for ensembles and beyond}
\label{chap:algos}
%% ======================================================================

\begin{objectives}
\guiding{How do we learn the weights of a forest, of a vote of boosted voters,
of a linear classifier or of a neural network by minimizing a certificate?}
This section applies the recipe of \cref{chap:selfbounding} to the ensemble
methods of the course, and then to linear classifiers, kernels and neural
networks. For each algorithm we specify the five ingredients of
\cref{algo:sb_template}: the voters and the organization of the data, the prior,
the objective, the optimizer and the certificate. The central tool for bagging is
the out-of-bag construction, which lets every example serve both to build some
voters and to certify the others. We run the algorithms on real data and look at
what they learn, including when they fail.
\end{objectives}

\subsection{Warm-up: a finite ensemble with an independent sample}

Before tackling bagging, where voters and bound share the same data, we start
with the simplest situation, the one of the decision stumps of \cref{fig:temperature}.
A set of $n$ voters is built on a sample $S_1$, for instance trees learned on
bootstrap samples of $S_1$, or stumps whose thresholds are chosen on $S_1$. A
second sample $S_2$ of size $m$ is used to compute the empirical risks
$\widehat R_{S_2}(h_t)$, and the prior is uniform over the $n$ voters, which is
legitimate since it only depends on $S_1$.

The remaining ingredients follow the toolbox of \cref{sec:sb_optim}. The objective can be the
PAC-Bayes-$\lambda$ bound, minimized by the Gibbs posterior
$\rho_t\propto e^{-\lambda m\widehat R_{S_2}(h_t)}$; the temperature $\lambda$ can be
scanned on a grid, since it only indexes posteriors, or optimized by alternating
with the closed-form $\lambda^\star$ (\cref{algo:fo} below). The certificate is
Seeger's bound on $S_2$ for the Gibbs classifier, or one of the bounds of
\cref{chap:mv} for the vote, the tandem losses being computed on $S_2$.

On \texttt{breast cancer}, with $n=300$ stumps and $m=120$ examples in $S_2$,
we obtained in \cref{chap:posteriors} a Seeger certificate of $0.256$ for a test
Gibbs risk of $0.114$. The certificate would be exactly the same if $\lambda$ had
been chosen by cross-validation, or by any other method: every choice of
$\lambda$ is certified. The drawback of this construction is that the examples of
$S_1$ only serve to build the voters, and those of $S_2$ only to weight and
certify them. Can every example play both roles? The bagging construction
below shows that it can.

\subsection{Bagging and random forests: out-of-bag PAC-Bayes}\label{sec:oob}

\subsubsection{The split construction}

The prior must not depend on the data used to evaluate the bound. With
bagging this looks problematic, since every tree is learned from $S$. Yet
each tree $h_t$ is learned on a bootstrap sample $S_t$ only, and the
\emph{out-of-bag} (OOB) examples $T_t=S\setminus S_t$ were not used to build it:
from the point of view of $h_t$, $T_t$ is a fresh test sample
(\cref{fig:oob}). The idea of \citet{thiemann2017strongly} is to use the
\emph{indices} of the trees as hypotheses, with a prior on the indices fixed in
advance, and to evaluate each tree on its own OOB sample.

\begin{figure}[t]
\centering
\begin{tikzpicture}[scale=0.55]
\foreach \t/\pat in {1/{1,1,0,1,0,1,1,0,1,1,0,1,1,1,0,1},2/{0,1,1,1,1,0,1,1,0,1,1,0,1,1,1,0},3/{1,0,1,1,1,1,0,1,1,0,1,1,0,1,1,1}}{
  \node[left] at (0,-\t*1.2+0.35) {\small $h_{\t}$};
  \foreach \b [count=\i] in \pat {
    \ifnum\b=1
      \fill[pbblue!45] (\i-1,-\t*1.2) rectangle (\i,-\t*1.2+0.7);
    \else
      \fill[pborange!70] (\i-1,-\t*1.2) rectangle (\i,-\t*1.2+0.7);
    \fi
    \draw[white] (\i-1,-\t*1.2) rectangle (\i,-\t*1.2+0.7);
  }
}
\node[left] at (0,0.35) {\small $S$};
\foreach \i in {1,...,16}{\draw (\i-1,0) rectangle (\i,0.7); \node[font=\tiny] at (\i-0.5,0.35) {$z_{\i}$};}
\fill[pbblue!45] (0,-4.6) rectangle (0.8,-4.1); \node[right,font=\small] at (0.9,-4.35) {in the bootstrap sample $S_t$ (training)};
\fill[pborange!70] (0,-5.5) rectangle (0.8,-5.0); \node[right,font=\small] at (0.9,-5.25) {out-of-bag $T_t=S\setminus S_t$ (evaluation)};
\node[right,font=\small,align=left] at (0,-6.2) {tandem loss of $(h_1,h_2)$: evaluated on $T_1\cap T_2$};
\end{tikzpicture}
\caption{Out-of-bag construction. Each tree is trained on its bootstrap
sample (blue) and its loss is estimated on its OOB sample (orange), which is
independent of the tree. The joint error of two trees is estimated on the
examples that are out-of-bag for both.}
\label{fig:oob}
\end{figure}

\begin{theorem}{PAC-Bayes with a split construction \citep{thiemann2017strongly,lorenzen2019pac}}{oob}
Let $h_1,\dots,h_n$ be voters, $h_t$ being learned on $S_t\subset S$ by any
procedure, and let $T_t=S\setminus S_t$ with $|T_t|\geq n_0$ for all $t$. Let
$\widehat L_t=\frac{1}{|T_t|}\sum_{(x,y)\in T_t}\ind{h_t(x)\neq y}$ be the OOB
loss of $h_t$ and $\pi$ a distribution on $\{1,\dots,n\}$ fixed in advance
(e.g.\ uniform). Then, with probability at least $1-\delta$ over $S\sim\D^m$ (and the
randomness of the bootstrap), for all $\rho\in\M(\{1,\dots,n\})$,
\begin{equation}\label{eq:oob}
\kl\Big(\sum_t\rho_t\widehat L_t\ \Big\|\ \sum_t\rho_t\RD(h_t)\Big)\leq\frac{\KL(\rho\|\pi)+\ln\frac{2\sqrt{n_0}}{\delta}}{n_0}.
\end{equation}
The same holds for the tandem losses $\widehat L_{tt'}$ computed on $T_t\cap T_{t'}$,
with $\KL(\rho^2\|\pi^2)=2\KL(\rho\|\pi)$ and $n_0\leq\min_{t,t'}|T_t\cap T_{t'}|$.
\end{theorem}

\begin{proofsketch}
We follow the proof of Theorem~\ref{thm:general}, the ``hypotheses'' being the
indices $t$, on which $\pi$ is a legitimate data-independent prior. Only the
third step changes: for each $t$, conditionally on $S_t$ the voter $h_t$ is fixed
and $T_t$ is an i.i.d.\ sample independent of it, so Maurer's lemma gives
$\E\,e^{|T_t|\kl(\widehat L_t\|\RD(h_t))}\leq2\sqrt{|T_t|}$. Since $|T_t|\geq n_0$,
Jensen's inequality yields $\E\,e^{n_0\kl(\widehat L_t\|\RD(h_t))}\leq(2\sqrt{|T_t|})^{n_0/|T_t|}\leq2\sqrt{n_0}$.
\end{proofsketch}

In plain words, each tree is graded only on examples it has never seen, and the
bound behaves as if we had an independent sample of size $n_0$; the posterior
over trees is then free, exactly as in \cref{chap:selfbounding}. The price is
that the certificate is computed with the size $n_0$ of the \emph{smallest} OOB set,
so it is important to make these sets large. With the standard bootstrap, each
OOB set contains about $(1-1/m)^m\approx e^{-1}\approx37\%$ of the sample, and the
intersection of two OOB sets about $e^{-2}\approx14\%$. Drawing bootstrap samples
of size $m/2$ instead of $m$ raises these proportions to about
$e^{-1/2}\approx61\%$ and $e^{-1}\approx37\%$, at the price of slightly weaker
trees; this is the choice made in the figures and notebooks of this course.

On the \texttt{digits} experiments below, the sizes are $m=1257$, $n_1=747$, $n_2=423$: the
first-order bound is computed with $n_1$ examples and the tandem bound with $n_2$
examples, instead of the $m$ training examples.

\subsubsection{Minimizing the first-order bound}\label{sec:algo_fo}

With the OOB losses in hand, the most natural first attempt is to minimize the
first-order bound. With the PAC-Bayes-$\lambda$ form \eqref{eq:lambda} of the bound, the optimal
posterior for a fixed $\lambda$ is the Gibbs posterior on the OOB losses
(Proposition~\ref{prop:gibbs_post}), and the optimal $\lambda$ for a fixed posterior
is given by \eqref{eq:lambda_star}. This suggests an alternating minimization,
in which each step solves one half of the problem exactly.

\begin{algorithm}[h]
\caption{PAC-Bayes-$\lambda$ minimization for a forest \citep{thiemann2017strongly}}\label{algo:fo}
\KwIn{OOB losses $\widehat L_1,\dots,\widehat L_n$, prior $\pi$, $n_0$, $\delta$}
$\lambda\leftarrow1$\;
\Repeat{convergence of $\lambda$}{
  $\rho_t\leftarrow\dfrac{\pi_t\,e^{-\lambda n_0\widehat L_t}}{\sum_{s}\pi_s\,e^{-\lambda n_0\widehat L_s}}$ for $t=1,\dots,n$\;
  $\lambda\leftarrow\dfrac{2}{\sqrt{2n_0\,\rho^\top\widehat L/\big(\KL(\rho\|\pi)+\ln\frac{2\sqrt{n_0}}{\delta}\big)+1}+1}$\;
}
\KwOut{$\rho$ and the certificate $2\,\klup\big(\rho^\top\widehat L,\frac{\KL(\rho\|\pi)+\ln(2\sqrt{n_0}/\delta)}{n_0}\big)$ on $\RD(B_\rho)$}
\end{algorithm}

\citet{thiemann2017strongly} proved that, under mild conditions, the
PAC-Bayes-$\lambda$ bound is strongly quasiconvex and that this procedure
converges to its global minimum. \Cref{tab:fo_iterations} follows it on a forest
of $100$ trees on \texttt{digits}. Starting from $\lambda=1$, the first Gibbs
posterior already has an empirical Gibbs risk of $0.128$ and a KL of $4.2$ nats;
$\lambda$ then drops to $0.37$ and is stable after three iterations, and the
bound on the Gibbs risk goes from $0.207$ to $0.206$. The algorithm is fast and
does exactly what it promises: it produces an excellent bound on the
\emph{Gibbs} risk.

\begin{table}[t]
\centering\small
\caption{Iterations of \cref{algo:fo} on a forest of $100$ trees on \texttt{digits}
($\delta=0.05$). On row $k$, $\rho$ is the Gibbs posterior for the value of $\lambda$
shown, and the PB-$\lambda$ bound is evaluated at the updated value of $\lambda$ (the one of
row $k+1$). The last column is the test risk of the majority vote $B_\rho$;
the uniform forest has a test risk of $0.044$. Later iterations are identical.}
\label{tab:fo_iterations}
\setlength{\tabcolsep}{4pt}
\begin{tabular}{rcccccc}
\toprule
iter. & $\lambda$ & $\rho^\top\widehat L$ & $\KL(\rho\|\pi)$ & PB-$\lambda$ bound on $R(G_\rho)$ & FO cert. & test $R(B_\rho)$ \\
\midrule
1 & 1.0000 & 0.1283 & 4.18 & 0.2070 & 0.3866 & 0.150 \\
2 & 0.3802 & 0.1293 & 3.76 & 0.2064 & 0.3863 & 0.150 \\
3 & 0.3735 & 0.1294 & 3.74 & 0.2064 & 0.3863 & 0.150 \\
4 & 0.3732 & 0.1294 & 3.74 & 0.2064 & 0.3863 & 0.150 \\
5 & 0.3732 & 0.1294 & 3.74 & 0.2064 & 0.3863 & 0.150 \\
\bottomrule
\end{tabular}
\end{table}

The last column of the table tells a different story: the test risk of the
vote is $0.150$, whereas the uniform forest had a test risk of $0.044$. As
observed by \citet{lorenzen2019pac} and \citet{masegosa2020second}, the vote
learned by minimizing a first-order bound is often much \emph{worse} than the
uniform vote. The Gibbs posterior concentrates on the few trees with the smallest
OOB loss; the left panel of \cref{fig:rf_weights} shows that it gives
essentially all its weight to two trees. These trees are good individually, but
a vote of two trees is not a forest, and the benefit of voting disappears. This
is a direct consequence of the first-order nature of the bound, which only sees
the average quality of the voters and ignores the correlations of their errors.
To protect the vote, the objective must see these correlations.

\subsubsection{Minimizing the tandem bound}

The tandem bound sees exactly what the first-order bound misses, since its
empirical term counts the errors made jointly by pairs of voters. Combined with
the PAC-Bayes-$\lambda$ relaxation, the tandem bound \eqref{eq:tnd_pb} gives the objective
\begin{equation}\label{eq:tnd_objective}
\mathcal B_{\text{TND}}(\rho,\lambda)=4\left[\frac{\rho^\top\widehat{\mathbf L}\rho}{1-\frac\lambda2}+\frac{2\KL(\rho\|\pi)+\ln\frac{2\sqrt{n_0}}{\delta}}{\lambda\big(1-\frac\lambda2\big)n_0}\right],
\qquad\widehat{\mathbf L}=(\widehat L_{tt'})_{t,t'},
\end{equation}
which is quadratic in $\rho$, plus the convex KL term. It has no closed-form
minimizer, but it is easily minimized by gradient descent on the softmax
parametrization $\rho=\mathrm{softmax}(\theta)$, alternated with the closed-form
update of $\lambda$ \citep{masegosa2020second}. The gradient with respect to $\rho$ is
\begin{equation}\label{eq:grad_tnd_lambda}
\nabla_\rho\mathcal B_{\text{TND}}=4\left[\frac{2\widehat{\mathbf L}\rho}{1-\frac\lambda2}+\frac{2\big(\ln\frac{\rho}{\pi}+\mathbf 1\big)}{\lambda\big(1-\frac\lambda2\big)n_0}\right],
\end{equation}
to be combined with \eqref{eq:softmax_grad}.

\begin{algorithm}[h]
\caption{Tandem bound minimization for a forest \citep{masegosa2020second}}\label{algo:tnd}
\KwIn{OOB tandem losses $\widehat{\mathbf L}$ (on OOB intersections), prior $\pi$, $n_0=\min_{t,t'}|T_t\cap T_{t'}|$, $\delta$}
$\theta\leftarrow\ln\pi$, $\lambda\leftarrow1$\;
\Repeat{convergence of $\lambda$ (or a maximal number of rounds)}{
  $\theta\leftarrow$ minimize $\mathcal B_{\text{TND}}(\mathrm{softmax}(\theta),\lambda)$ with L-BFGS, using \eqref{eq:grad_tnd_lambda} and \eqref{eq:softmax_grad}\;
  $\rho\leftarrow\mathrm{softmax}(\theta)$\;
  $\lambda\leftarrow\dfrac{2}{\sqrt{2n_0\,\rho^\top\widehat{\mathbf L}\rho/\big(2\KL(\rho\|\pi)+\ln\frac{2\sqrt{n_0}}{\delta}\big)+1}+1}$\;
}
\KwOut{$\rho$ and the certificate $4\,\klup\big(\rho^\top\widehat{\mathbf L}\rho,\frac{2\KL(\rho\|\pi)+\ln(2\sqrt{n_0}/\delta)}{n_0}\big)$ on $\RD(B_\rho)$}
\end{algorithm}

In words, the quadratic term $\rho^\top\widehat{\mathbf L}\rho$ penalizes pairs of voters that
err \emph{together}, so the optimal posterior has no reason to concentrate on
a few similar trees: diversity is rewarded instead of ignored. Alternatively, the kl-form certificate itself can be
minimized with the gradient \eqref{eq:grad_tnd_cert} of \cref{sec:sb_optim}; the
dynamics of this approach were shown in \cref{fig:sb_dynamics}.

\subsubsection{Comparing the learned posteriors}

We can now put the two algorithms side by side, together with the uniform
weights and a C-bound minimizer. \Cref{tab:rf_posteriors} and \cref{fig:rf_weights} compare, on the same forest,
four posteriors: the uniform one, the minimizers of the first-order and of the
tandem bounds, and the minimizer of the PAC-Bayesian C-bound \eqref{eq:cbound_pb},
obtained by a direct numerical minimization. The first-order posterior has the
smallest Gibbs risk and the smallest first-order certificate, but it keeps two
effective trees and its vote is three times worse than the others.

The tandem posterior keeps an effective number of $58$ trees, spread over the good and
the average trees (right panel of \cref{fig:rf_weights}): its vote is as good as
the uniform one ($0.043$), and it has the smallest certificate of all, $0.361$.
The C-bound posterior is in between, with $31$ effective trees and a vote of
comparable quality, but its certificate ($0.531$) is looser than the tandem
certificate, because it needs two kl inversions and a ratio. This illustrates the
remark following Proposition~\ref{prop:compare}: the C-bound is the best
\emph{oracle} bound, but not necessarily the best \emph{certificate}.

\begin{table}[t]
\centering\small
\caption{Four posteriors on the same forest of $100$ trees on \texttt{digits}:
test risks of the vote and of the Gibbs classifier, the three OOB certificates on the risk
of the vote ($\delta=0.05$), and the effective number of trees $1/\sum_t\rho_t^2$.}
\label{tab:rf_posteriors}
\setlength{\tabcolsep}{3.5pt}
\begin{tabular}{lcccccc}
\toprule
posterior & $R(B_\rho)$ & $R(G_\rho)$ & FO cert. & TND cert. & C-bound cert. & eff.\ trees \\
\midrule
uniform & 0.044 & 0.190 & 0.460 & 0.387 & 0.559 & 100.0 \\
FO-optimized & 0.150 & 0.158 & 0.386 & 0.666 & 0.596 & 2.0 \\
TND-optimized & 0.043 & 0.183 & 0.436 & 0.361 & 0.538 & 58.1 \\
C-bound-optimized & 0.048 & 0.176 & 0.419 & 0.369 & 0.531 & 30.8 \\
\bottomrule
\end{tabular}
\end{table}

\begin{figure}[t]
\centering
\includegraphics[width=\linewidth]{rf_weights}
\caption{Weights learned on a forest of 100 trees on \texttt{digits}. Left:
sorted weights of the four posteriors (log scale). Right: weight of each tree
as a function of its OOB loss for the first-order and tandem posteriors: the
first-order posterior is a hard selection of the two best trees, the tandem
posterior a soft one.}
\label{fig:rf_weights}
\end{figure}

Is this behaviour specific to \texttt{digits}? \Cref{fig:rf_opt} repeats the
comparison between the uniform, first-order and tandem posteriors on four
datasets. The picture is consistent: minimizing the
first-order bound always improves the first-order certificate, but it degrades
the vote on all four datasets, sometimes severely (from $0.044$ to
$0.150$ on \texttt{digits}, from $0.074$ to $0.184$ on the simulated data).
Minimizing the tandem bound never degrades the vote and always yields the best
tandem certificate. The certificates remain far from the test risks --- this is
the price of certifying a vote with a second-order Markov inequality and a
doubled KL --- but they are valid and obtained without any test data.

\begin{figure}[t]
\centering
\includegraphics[width=\linewidth]{rf_optimisation}
\caption{Forests of 100 trees (bootstrap samples of size $m/2$), on four
datasets. For the uniform posterior, the posterior minimizing the first-order
(FO) bound, and the posterior minimizing the tandem (TND) bound, we report the
test risk of the majority vote and the two OOB certificates ($\delta=0.05$).
Generated with \texttt{pbtools.py} (notebook 4).}
\label{fig:rf_opt}
\end{figure}

\subsection{Boosting-like algorithms: learning the weights of a vote}

Bagging builds its voters independently and only asks us to weight them.
Boosting builds them sequentially and learns their weights at the same time;
we now see how the PAC-Bayesian viewpoint, and in particular the C-bound,
shapes algorithms of this kind.

\subsubsection{AdaBoost from the PAC-Bayesian viewpoint}

AdaBoost greedily minimizes the exponential loss $\frac1m\sum_ie^{-y_if(x_i)}$
of $f=\sum_t\alpha_th_t$, an upper bound on the training error of the vote
\citep{freund1997decision}. In PAC-Bayesian terms, it learns a posterior
$Q\propto(\alpha_t)_t$ whose margins $M_Q$ on $S$ are large, the normalized
margins being $y_if(x_i)/\|\alpha\|_1$; this is the explanation of
\citet{schapire1998boosting}, illustrated in \cref{fig:adaboost}. But nothing in
AdaBoost controls the divergence $\KL(Q\|P)$ or the disagreement of the voters.
The algorithms of this subsection instead learn the weights of a set of voters,
possibly produced by boosting, by minimizing a C-bound.

\subsubsection{MinCq: minimizing the C-bound}

The C-bound rewards votes whose margin is large on average and varies little.
Recall that $\mathcal C_Q=1-\frac{(\E M_Q)^2}{\E M_Q^2}$. For a fixed value $\mu$ of the
first moment of the margin, minimizing the C-bound amounts to minimizing its
second moment. \citet{laviolette2011from} and \citet{germain2015risk} proposed
the algorithm
\begin{equation}\label{eq:mincq}
\textbf{MinCq:}\qquad \min_{Q}\ \widehat{\E}_S\big[M_Q^2\big]\quad\text{s.t.}\quad\widehat{\E}_S\big[M_Q\big]=\mu,\quad Q\text{ quasi-uniform},
\end{equation}
where $\Hcal$ is \emph{self-complemented} ($h\in\Hcal\Rightarrow-h\in\Hcal$) and a
\emph{quasi-uniform} posterior satisfies $Q(h)+Q(-h)=\frac1{|\Hcal|/2}$. Since
$M_Q$ is linear in $Q$, $\widehat\E_S[M_Q^2]$ is a quadratic form in the weights
and \eqref{eq:mincq} is a quadratic program.

The restriction to quasi-uniform posteriors is the key trick. On the one hand, it does not reduce the set of
votes that can be expressed, since any vote can be represented by a
quasi-uniform posterior with the same predictions. On the other hand,
PAC-Bayesian bounds hold for quasi-uniform posteriors \emph{without the KL term}
\citep{germain2015risk}, so that the empirical C-bound itself is controlled.

The hyperparameter $\mu$, the targeted first moment of the margin, plays the role of
a regularization parameter and is selected by cross-validation. MinCq has been
extended to incorporate prior knowledge on the weights
\citep[P-MinCq,][]{bellet2014learning} and to a column-generation algorithm that
adds voters one at a time \citep[CqBoost,][]{roy2016column}.

\subsubsection{CB-Boost: greedy C-bound minimization}

MinCq weights a fixed set of voters. To grow the vote one voter at a time, as
AdaBoost does, \citet{bauvin2020fast} proposed CB-Boost, a boosting-like algorithm that builds
the vote greedily by minimizing the empirical C-bound. At each round, it adds a
voter and chooses its weight so as to decrease the empirical C-bound; since this
quantity is a ratio of two quadratic functions of the weights, the optimal weight
of the new voter has a closed form, which makes each round as cheap as a round of
AdaBoost. The empirical C-bound decreases at each iteration, and the final vote
can be certified with a PAC-Bayesian C-bound.

\subsubsection{Direct minimization of PAC-Bayesian C-bounds}

MinCq and CB-Boost minimize the \emph{empirical} C-bound and certify afterwards.
A fully self-bounding alternative is to minimize the certificate itself:
\citet{viallard2021self} proposed to minimize by (stochastic) gradient descent
the \emph{PAC-Bayesian} C-bounds themselves, that is the right-hand side of
\eqref{eq:cbound_pb} and tighter joint versions, over the posterior
$\rho=\mathrm{softmax}(\theta)$. The C-bound \eqref{eq:cbound_pb} is a composition of
$\klup$ (for the Gibbs risk), $\kllow$ (for the disagreement) and a smooth
function, so its gradient follows from Proposition~\ref{prop:klgrad} and the chain
rule: with $\bar r=\klup(\rho^\top\widehat L,\varepsilon_1)$, $\underline d=\kllow(\rho^\top\widehat D\rho,\varepsilon_2)$ and
$\mathcal C=1-\frac{(1-2\bar r)^2}{1-2\underline d}$,
\[
\frac{\partial\mathcal C}{\partial\bar r}=\frac{4(1-2\bar r)}{1-2\underline d},\qquad
\frac{\partial\mathcal C}{\partial\underline d}=-\frac{2(1-2\bar r)^2}{(1-2\underline d)^2}.
\]
The learned vote comes with its certificate by construction.

\citet{zantedeschi2021learning}
go one step further and consider \emph{stochastic majority votes}, whose weights
are themselves drawn from a Dirichlet posterior. For such votes the expected risk
can be computed, or tightly approximated, and bounded directly, which avoids the
loose Gibbs-to-vote relations of \cref{chap:mv}.

\subsection{Linear classifiers and kernel methods}\label{sec:linear}

The recipe is not limited to finite ensembles. Linear classifiers are the
natural next step: they are the one case where every term of the bound can be
written explicitly.

\subsubsection{Gaussian posteriors on linear classifiers}

Consider linear voters $h_{\mathbf v}(x)=\sign(\langle\mathbf v,x\rangle)$,
$\mathbf v\in\R^d$, a prior $P=\mathcal N(\mathbf 0,I_d)$ and posteriors
$Q_{\mathbf w}=\mathcal N(\mathbf w,I_d)$. By symmetry of the Gaussian around
$\mathbf w$, the majority vote is
\[
B_{Q_{\mathbf w}}(x)=\sign\big(\E_{\mathbf v}\sign\langle\mathbf v,x\rangle\big)=\sign\langle\mathbf w,x\rangle:
\]
the vote of an infinite number of linear classifiers is the linear classifier
$\mathbf w$ itself. For $\mathbf v\sim Q_{\mathbf w}$, the random variable
$y\langle\mathbf v,x\rangle$ is Gaussian with mean $y\langle\mathbf w,x\rangle$ and
variance $\|x\|^2$, so the empirical Gibbs risk is
\begin{equation}\label{eq:probit}
\RS(G_{Q_{\mathbf w}})=\frac1m\sum_{i=1}^m\Phi\Big(-\frac{y_i\langle\mathbf w,x_i\rangle}{\|x_i\|}\Big),
\end{equation}
where $\Phi$ is the standard Gaussian cumulative distribution function: the 0-1
loss is replaced by a smooth \emph{probit} loss of the normalized margin.
Finally, the complexity is $\KL(Q_{\mathbf w}\|P)=\frac12\|\mathbf w\|^2$.
\Cref{fig:linear} shows the resulting trade-off. Multiplying $\mathbf w$ by a
constant does not change the vote, but it makes the sampled classifiers agree
more with it: the Gibbs risk decreases towards the risk of the vote, while the
KL grows quadratically.

\begin{figure}[t]
\centering
\includegraphics[width=0.85\linewidth]{linear_gaussian}
\caption{Gaussian posterior $\mathcal N(\mathbf w,I)$ over linear classifiers:
40 sampled hyperplanes (grey) and the majority vote $\sign\langle\mathbf w,x\rangle$
(black). Multiplying $\mathbf w$ by $5$ does not change the vote but makes the
sampled voters agree with it: the Gibbs risk decreases while the KL grows
quadratically.}
\label{fig:linear}
\end{figure}

Plugging these quantities into Catoni's bound \eqref{eq:catoni} and removing
the monotone transformation gives the objective of the algorithm PBGD
\citep{germain2009pac}:
\begin{equation}\label{eq:pbgd}
\min_{\mathbf w\in\R^d}\ F(\mathbf w)=C\sum_{i=1}^m\Phi\Big(-\frac{y_i\langle\mathbf w,x_i\rangle}{\|x_i\|}\Big)+\frac12\|\mathbf w\|^2.
\end{equation}
It is instructive to compare it with the soft-margin SVM,
$\min_{\mathbf w}C\sum_i\max(0,1-y_i\langle\mathbf w,x_i\rangle)+\frac12\|\mathbf w\|^2$.
PAC-Bayes \emph{explains} the SVM objective: the regularization is the KL to the
prior, the loss is a surrogate of the Gibbs risk, and the norm of $\mathbf w$
plays the role of the inverse of a margin \citep{langford2002pac,mcallester2003simplified}.
The two objectives differ mainly in the loss, and this difference has a
practical consequence. The objective \eqref{eq:pbgd} is not convex; its gradient is
\[
\nabla F(\mathbf w)=-C\sum_{i=1}^m\varphi\Big(\frac{y_i\langle\mathbf w,x_i\rangle}{\|x_i\|}\Big)\frac{y_ix_i}{\|x_i\|}+\mathbf w,
\]
where $\varphi$ is the standard Gaussian density. Each example pushes $\mathbf w$
in its direction, with a weight that is maximal for points close to the decision
boundary and vanishes for points with a large margin, positive or negative:
unlike the hinge loss, badly misclassified points are eventually ignored, which
makes PBGD robust to label noise. It is optimized by gradient descent with
several restarts, and has a kernelized version \citep{germain2009pac}.

\begin{example}[PBGD selects its own regularization]
As a self-bounding algorithm, PBGD has an important advantage: the parameter $C$
only indexes posteriors, so it can be chosen by minimizing the certificate
itself, without cross-validation. \Cref{fig:pbgd} runs PBGD on a two-dimensional
problem with $m=200$ examples and $8\%$ of flipped labels, for $25$ values of $C$
between $10^{-2}$ and $10^2$, and computes Seeger's bound for each solution. For
small $C$, $\mathbf w$ is short ($\|\mathbf w\|=0.41$ for $C=0.01$), the sampled
classifiers are almost random and the certificate is $0.54$. For large $C$, the
Gibbs risk approaches the training error of the vote but the KL explodes
($\|\mathbf w\|=13.6$ for $C=100$), and the certificate climbs back to $0.58$. The
best certificate, $0.28$, is reached for $C=1$, with $\|\mathbf w\|=3.5$, an
empirical Gibbs risk of $0.133$ and a training error of the vote of $0.110$. The
right panel shows the hyperplanes learned with a small $C$, with the $C$ selected
by the certificate, and with the largest $C$; no validation sample was used to
choose among them.
\end{example}

\begin{figure}[t]
\centering
\includegraphics[width=\linewidth]{pbgd}
\caption{PBGD on a two-dimensional problem ($m=200$, $8\%$ label noise). Left:
empirical Gibbs risk, complexity $\|\mathbf w\|^2/(2m)$ and Seeger certificate as
functions of $C$. Right: hyperplanes learned for three values of $C$.}
\label{fig:pbgd}
\end{figure}

Data-dependent Gaussian priors, centred on an SVM learned on part of the data,
give some of the tightest bounds for SVMs \citep{ambroladze2006tighter,parrado2012pac}.
The prior need not be Gaussian, however: a kernel can itself be read as a prior
over voters.

\subsubsection{Kernels as priors: random Fourier features}

For a shift-invariant kernel $k(x,x')=k(x-x')$, Bochner's theorem gives
$k(x-x')=\E_{\omega\sim p}\cos\big(\omega^\top(x-x')\big)$, where $p$ is the Fourier
transform of $k$ \citep{rahimi2007random}. \citet{letarte2019pseudo} interpret this
identity as a \emph{majority vote}: the voters are the functions
$h_\omega(x,x')=\cos(\omega^\top(x-x'))$ and the kernel is their vote under the
\emph{prior} $p$. Learning a kernel adapted to the task then amounts to learning
a \emph{posterior} over frequencies. With $N$ frequencies $\omega_1,\dots,\omega_N$
drawn from $p$, the pseudo-posterior $Q(\omega_j)\propto e^{-\beta\,\widehat L_S(h_{\omega_j})}$
is exactly a Gibbs posterior, and a PAC-Bayesian bound certifies the learned
kernel. This is closely related to the boosting-with-random-Fourier-features
algorithm studied in the first project of the course.

\begin{example}[Random Fourier voters on the two moons]
A simpler variant uses the random features directly as voters,
$h_{\omega,b}(x)=\sign\cos(\omega^\top x+b)$, with $\omega\sim\mathcal N(0,2\gamma I)$ and
$b\sim\mathcal U[0,2\pi]$, together with their opposites. \Cref{fig:rff} shows one
such voter on the two moons dataset, and the vote of $4000$ of them under Gibbs
posteriors at three temperatures, learned on $900$ examples. Under the uniform
prior the vote is undecided, since each voter comes with its opposite. A very
small tilt of the prior ($\lambda=0.001$, $\KL=0.005$ nats) already gives a smooth
decision boundary with a test risk of $0.067$: the vote then behaves like a
kernel classifier. Larger values of $\lambda$ (lower temperatures) concentrate the posterior on fewer
frequencies and give more irregular boundaries, with test risks of $0.11$
($\lambda=0.02$) and $0.13$ ($\lambda=0.5$). Note that the best vote corresponds to a
posterior whose empirical Gibbs risk is $0.49$: the Gibbs certificate is
vacuous, while the vote is excellent. This is an extreme instance of the gap
studied in \cref{chap:mv}, and a reminder that votes must be certified with
second-order bounds.
\end{example}

\begin{figure}[t]
\centering
\includegraphics[width=\linewidth]{rff_votes}
\caption{Random Fourier voters on the two moons. Left: one voter. Right: the
majority vote under Gibbs posteriors at three temperatures (test risks in the titles).}
\label{fig:rff}
\end{figure}

\subsection{Beyond ensembles: stochastic neural networks and Gaussian processes}

The Gaussian posteriors used for linear classifiers carry over to models with
many more continuous parameters, and the same recipe applies. For a neural
network with weights $\mathbf w\in\R^D$, consider a Gaussian posterior
$Q=\mathcal N(\boldsymbol\mu,\mathrm{diag}(\boldsymbol\sigma^2))$ and a Gaussian prior
$P=\mathcal N(\boldsymbol\mu_0,\sigma_0^2I)$; drawing a network from $Q$ at each prediction
defines a \emph{stochastic neural network}. The centre $\boldsymbol\mu_0$ of the prior is
the random initialization of the network or, much better, a network trained on a
part $S_1$ of the data (Proposition~\ref{prop:ddprior}); the scale $\sigma_0$ is chosen
on a grid with a union bound. As for linear classifiers, the KL divergence has the closed form
\[
\KL(Q\|P)=\frac12\sum_{j=1}^D\Big(\frac{\sigma_j^2+(\mu_j-\mu_{0,j})^2}{\sigma_0^2}-1+\ln\frac{\sigma_0^2}{\sigma_j^2}\Big).
\]
The objective is one of the relaxations of Seeger's bound of \cref{sec:sb_optim},
in which the 0-1 loss is replaced by a bounded surrogate, for instance the
cross-entropy clipped to $[0,\ell_{\max}]$ and rescaled to $[0,1]$, and it is
minimized by stochastic gradient descent on $(\boldsymbol\mu,\ln\boldsymbol\sigma)$ with the
reparametrization trick. The certificate, finally, uses the 0-1 loss: the
empirical Gibbs risk is estimated with $N$ networks sampled from $Q$, bounded by
a first kl inversion, and plugged into Seeger's bound, as in \cref{fig:montecarlo}.

\citet{dziugaite2017computing} obtained in this way the first \emph{non-vacuous}
bounds for over-parametrized networks on MNIST, and \citet{perez2021tighter}
showed that such PAC-Bayesian training yields networks with competitive accuracy
together with tight certificates. \citet{reeb2018learning} applied the same
principle to learn the hyperparameters of Gaussian processes by minimizing a
PAC-Bayesian bound, and \citet{lotfi2022pac} obtained non-vacuous bounds for
large models with compression-based priors. Finally, \emph{deep ensembles}
\citep{lakshminarayanan2017simple} are majority votes of networks, to which the
second-order bounds of \cref{chap:mv} and the algorithms of this section apply.

\subsection{The running example, one last time}

Before summarizing the algorithms, let us step back and look at the whole
journey of our running example.

\begin{running}
Let us gather what we have learned about our forest of $100$ trees on
\texttt{digits}. Its vote errs on $4.4\%$ of the test digits, while a tree drawn
at random errs on $19\%$ of them; \cref{chap:mv} explained the gap by the diversity
of the trees, whose disagreement ($0.27$) exceeds their Gibbs risk. Setting aside
the $540$ test digits certified the vote at $0.069$ (\cref{chap:tools}). Without
any test set, the out-of-bag tandem bound certified the uniform vote at $0.387$
(\cref{chap:mv}), and learning the weights by minimizing this bound lowered the
certificate to $0.361$ while keeping the vote as accurate ($0.043$). Minimizing
the first-order bound instead gave the best certificate on the Gibbs risk, but
destroyed the vote ($0.150$). Three lessons summarize the course: the diversity
of the voters is what makes a vote work; a certificate for a vote must see this
diversity, which requires second-order quantities; and a certificate that sees it
can safely be used as a learning objective.
\end{running}

\subsection{Summary}

The following table summarizes the algorithms of this section with the five
ingredients of the self-bounding recipe.

\begin{center}
\small
\renewcommand{\arraystretch}{1.3}
\begin{tabular}{p{2.6cm}p{2.1cm}p{2.9cm}p{2.1cm}p{1.9cm}}
\toprule
Algorithm & Posterior & Objective & Optimizer & Certifies \\
\midrule
Gibbs posterior & finite, any & linear bound & closed form & Gibbs, vote \\
PB-$\lambda$ \citep{thiemann2017strongly} & finite (OOB) & PAC-Bayes-$\lambda$ & alternating & Gibbs \\
TND \citep{masegosa2020second} & finite (OOB) & tandem-$\lambda$ & L-BFGS + $\lambda$ & vote \\
MinCq \citep{laviolette2011from} & quasi-uniform & empirical C-bound & QP & vote \\
CB-Boost \citep{bauvin2020fast} & greedy & empirical C-bound & closed-form steps & vote \\
Self-bounding C-bound \citep{viallard2021self} & finite & PAC-Bayesian C-bound & gradient (kl inv.) & vote \\
StocMV \citep{zantedeschi2021learning} & Dirichlet & bound on stochastic vote & gradient & stochastic vote \\
PBGD \citep{germain2009pac} & Gaussian, linear & Catoni & gradient & Gibbs, vote \\
PBB \citep{perez2021tighter} & Gaussian, NN & Seeger relaxations & SGD, reparam. & Gibbs \\
\bottomrule
\end{tabular}
\end{center}

\begin{keypoints}
To certify an ensemble, each voter must be evaluated on examples it has not
seen: an independent sample, or its out-of-bag examples (Theorem~\ref{thm:oob}).
Minimizing a first-order bound gives an excellent bound on the Gibbs risk but
concentrates the weights on a few voters and destroys the diversity of the vote;
second-order objectives such as the tandem bound keep it, and give the best
certificates for votes. The C-bound leads to a family of algorithms (MinCq,
CB-Boost, direct minimization) which explicitly trade the strength of the voters
against their disagreement. The same recipe gives PAC-Bayesian versions of SVMs
(PBGD, which selects its own regularization), of kernel learning and of
stochastic neural networks, where data-dependent priors are the key ingredient.
\end{keypoints}

\begin{quickchecks}
\item Why can each tree of a bagged forest be evaluated on its out-of-bag examples without breaking the bound?
\item Why does minimizing the first-order bound degrade the vote?
\item What does the parameter $\mu$ of MinCq control?
\item In PBGD, which quantity plays the role of the regularization of the SVM?
\end{quickchecks}

\begin{exercises}
\begin{exercise}
Derive the gradient \eqref{eq:grad_tnd_lambda} of the tandem objective, and check
it numerically with finite differences (notebook~4).
\end{exercise}
\begin{exercise}
Explain why the uniform prior over the trees of a random forest would \emph{not} be
a valid prior if the empirical losses were computed on the whole training sample.
\end{exercise}
\begin{exercise}
In PBGD, show that when all the training points have a large normalized margin,
the gradient of the loss term vanishes, and interpret the resulting value of $\|\mathbf w\|$.
\end{exercise}
\begin{exercise}
With notebook~4, compute the test disagreement of the C-bound posterior of
\cref{tab:rf_posteriors} and check whether the condition $d\geq R$ of
Proposition~\ref{prop:compare} holds. Why is its certificate nevertheless
larger than the tandem certificate?
\end{exercise}
\begin{exercise}
$(\star)$ For a stochastic network with $D=10^5$ weights, $\sigma_j=\sigma_0=0.03$ and
$\|\boldsymbol\mu-\boldsymbol\mu_0\|=3$, compute $\KL(Q\|P)$. How many examples are needed for the
complexity term $\KL/m$ to be below $0.01$? What does this suggest about data-dependent priors?
\end{exercise}
\begin{exercise}
\textbf{(Comprehension)} True or false? Justify briefly. (a)~The OOB bound is invalid because
every training example has been used to build some of the trees. (b)~Adding trees to a
forest with the uniform posterior increases $\KL(\rho\|\pi)$. (c)~In PBGD, the parameter $C$
must be selected by cross-validation, since the bound cannot be used to choose it.
(d)~The first-order posterior of \cref{tab:rf_posteriors} has a small Gibbs risk, hence
necessarily a good vote.
\end{exercise}
\begin{exercise}
\textbf{(Application)} For the \texttt{digits} forest ($m=1257$), compute the expected
fraction and number of examples that are out-of-bag for one tree, and for two given trees,
when the bootstrap samples have size $m$ and when they have size $\lfloor m/2\rfloor=628$.
Compare with the values $n_1=747$ and $n_2=423$ used in the notes, and explain the
difference.
\end{exercise}
\begin{exercise}
\textbf{(Application)} Run \cref{algo:fo} by hand (with a calculator or \texttt{pbtools}) on
$n=4$ trees with OOB losses $\widehat L=(0.15,0.17,0.20,0.25)$, uniform prior, $n_0=700$
and $\delta=0.05$: compute $\rho$ for $\lambda=1$, the updated $\lambda$, and $\rho$ again.
Compare the final first-order certificate with the one of the uniform posterior, and
comment on the weights.
\end{exercise}
\begin{exercise}
\textbf{(Application)} For a Gaussian posterior $Q_{\mathbf w}=\mathcal N(\mathbf w,I_d)$ with
$\|\mathbf w\|=2$, three training points have normalized margins
$y_i\langle\mathbf w,x_i\rangle/\|x_i\|$ equal to $2$, $0.5$ and $-1$. Compute the empirical
Gibbs risk \eqref{eq:probit}, the training error of the vote and $\KL(Q_{\mathbf w}\|P)$.
Redo the computation for $3\mathbf w$, and relate the result to \cref{fig:linear}.
\end{exercise}
\begin{exercise}
\textbf{(Proof)} Let $Q_{\mathbf w}=\mathcal N(\mathbf w,I_d)$ over linear voters
$h_{\mathbf v}(x)=\sign\langle\mathbf v,x\rangle$, and $x\neq0$. Show that
$\P_{\mathbf v\sim Q_{\mathbf w}}\big(y\langle\mathbf v,x\rangle\leq0\big)=\Phi\big(-y\langle\mathbf w,x\rangle/\|x\|\big)$,
and deduce that the majority vote is $B_{Q_{\mathbf w}}(x)=\sign\langle\mathbf w,x\rangle$ (up to the
tie $\langle\mathbf w,x\rangle=0$) and the formula \eqref{eq:probit}.
\end{exercise}
\end{exercises}
